\section{Frame Identifiability and Sequential Amplification}
\label{sec:theory}

We ask whether labels can compensate for a missing frame in any response-only policy and whether a constant refinement can redirect a current selector.

\paragraph{Response-compatible worlds.}
Let $O(E)=\{(e,a_e):e\in E\}$ be the public response view. Frames exposing the same $O(E)$ and label transcript are indistinguishable to a frame-blind policy: labels reveal behavior, not whether entries are independent roots. A fresh hidden-frame choice at every step makes this ambiguity accumulate.

Let $V^\pi_{\Phi,B}$ be the expected information acquired by policy $\pi$ about $B$ latent binary candidate coordinates, and $V^*_{\Phi,B}$ the optimum with known $\Phi$. Define $I_{\mathrm{hi}}=0.540852$ and $I_{\mathrm{lo}}=0.459148$ bits.

\begin{theorem}[Tight cost of a missing evidence frame]
\label{thm:sequential}
For every integer $B\geq1$, there is one fixed public response registry, a pool with two binary queries per coordinate, and $2^B$ response-compatible frames such that every randomized frame-blind policy has, on some compatible frame,
\[
V^*_{\Phi,B}-V^\pi_{\Phi,B}
\geq \frac{B}{2}(I_{\mathrm{hi}}-I_{\mathrm{lo}})
=0.040852B\ \text{bits}.
\]
The bound is tight; the corresponding tight minimax excess coordinate-wise Hamming risk is $B/24$.
\end{theorem}

One latent frame bit per coordinate determines which query polarity yields $I_{\mathrm{hi}}$. Before that coordinate is queried, the bit is independent of prior observations, so a frame-blind policy chooses the informative polarity with probability at most one half. Revisiting cannot improve the minimax value. Averaging gives a fixed hard frame and independent randomization attains equality (Appendix~\ref{app:proofs}). The exponential construction proves necessity, not typical registry size.

\begin{theorem}[Constant-registry full-path amplification]
\label{thm:amplification}
For every odd $B\geq3$ and finite Select-LLM temperature $\tau>0$, there is a pool with five canonical behaviors and $2B$ items such that adding four exact aliases produces nine entries, redirects all $B$ acquisitions to a disjoint block, and raises cumulative regret from zero to at least
\[
\frac{B^2-1}{4B}=\Omega(B),
\]
with terminal regret $(B-1)/(2B)$. The two best canonical behaviors and all existing responses are unchanged.
\end{theorem}

Four aliases move one behavior's mass above one half along an alternating-label path, making every bad query preferable and reinforcing the condition after each selection. Counting odd steps gives both regret bounds (Appendix~\ref{app:proofs}). This is mechanism-specific, and exact hashing blocks the literal construction; Section~\ref{sec:protocol} therefore tests response-distinct variants.

\begin{proposition}[Selector-agnostic root factorization]
\label{prop:factorization}
Suppose $\rho$ is authenticated, each root has fixed mass $w_r$, and a selector's state, action kernel, and decision depend on entries only through root summaries $S_r$. If $S_r$ is invariant to an allowed within-root refinement, then the complete query transcript and root decisions are invariant in distribution.
\end{proposition}

Transcript induction proves the claim (Appendix~\ref{app:proofs}). It covers uncertainty, diversity, and mutual-information rules only when their implemented state actually factorizes this way. Exact quotienting handles literal response copies, not response-distinct variants, and predictions cannot authenticate $\rho$. Appendix~\ref{app:v21-stability} gives the approximate counterpart: if the root-level next-step kernels differ by at most $\varepsilon_t$ in total variation, complete transcript distance is at most $1-\prod_t(1-\varepsilon_t)\leq\sum_t\varepsilon_t$, with the same factor controlling any bounded selection cost.
